\documentclass[../../main.tex]{subfiles}

\begin{document}

\section{Riemann 可积与 Lebesgue 可积的刻画}\label{section:ch05-Riemann可积刻画}

\begin{lemma}\label{lemma:ch05-夹逼}
令 $\{\varphi_{n}\}$ 和 $\{\psi_{n}\}$ 为函数序列, 它们在 $E$ 上可积, 使得 $\varphi_{n}$ 在 $E$ 上递增而 $\psi_{n}$ 在 $E$ 上递减. 令 $E$ 上的函数 $f$ 具有性质: 在 $E$ 上对所有 $n$,
\begin{align*}
\varphi_{n}\leqslant f\leqslant\psi_{n}.
\end{align*}
若
\begin{align*}
\lim_{n\to\infty}\int_{E}[\psi_{n}-\varphi_{n}]=0,
\end{align*}
则在 $E$ 上 a.e. 逐点地 $\{\varphi_{n}\}\to f$, 逐点地 $\{\psi_{n}\}\to f$, $f$ 在 $E$ 上可积,
\begin{align*}
\lim_{n\to\infty}\int_{E}\varphi_{n}=\int_{E}f\quad\text{且}\quad\lim_{n\to\infty}\int_{E}\psi_{n}=\int_{E}f.
\end{align*}
\end{lemma}

\begin{proof}
对 $E$ 中的 $x$, 定义
\begin{align*}
\varphi^{*}(x)=\lim_{n\to\infty}\varphi_{n}(x)\quad\text{与}\quad\psi^{*}(x)=\lim_{n\to\infty}\psi_{n}(x).
\end{align*}
函数 $\varphi^{*}$ 和 $\psi^{*}$ 是正确定义的, 这是由于单调扩充实数列收敛到一个扩充实数, 且它们是可测的, 这是由于它们都是可测函数序列的逐点极限. 我们有不等式
\begin{align}
\text{在 }E\text{ 上对所有 }n,\quad\varphi_{n}\leqslant\varphi^{*}\leqslant f\leqslant\psi^{*}\leqslant\psi_{n}.\label{align:ch05-夹逼不等式}
\end{align}
根据非负可测函数的积分的单调性与线性, 对所有 $n$,
\begin{align*}
0\leqslant\int_{E}(\psi^{*}-\varphi^{*})\leqslant\int_{E}(\psi_{n}-\varphi_{n}).
\end{align*}
因此
\begin{align*}
0\leqslant\int_{E}(\psi^{*}-\varphi^{*})\leqslant\lim_{n\to\infty}\int_{E}(\psi_{n}-\varphi_{n})=0.
\end{align*}
由于 $\psi^{*}-\varphi^{*}$ 是一个非负可测函数而 $\int_{E}(\psi^{*}-\varphi^{*})=0$, 第 4 章的命题\cref{proposition:ch04-积分为零}告诉我们在 $E$ 上 a.e. $\psi^{*}=\varphi^{*}$. 但在 $E$ 上 $\varphi^{*}\leqslant f\leqslant\psi^{*}$. 因此
\begin{align*}
\text{在 }E\text{ 上 a.e. 逐点地 }\{\varphi_{n}\}\to f\text{ 且 }\{\psi_{n}\}\to f.
\end{align*}
因此 $f$ 是可测的. 观察到由于在 $E$ 上 $0\leqslant f-\varphi_{1}\leqslant\psi_{1}-\varphi_{1}$ 且 $\psi_{1}$ 和 $\varphi_{1}$ 在 $E$ 上是可积的, 我们从积分的比较判别法推出 $f$ 在 $E$ 上是可积的. 从不等式\eqref{align:ch05-夹逼不等式}推出对所有 $n$,
\begin{align*}
0\leqslant\int_{E}\psi_{n}-\int_{E}f=\int_{E}(\psi_{n}-f)\leqslant\int_{E}(\psi_{n}-\varphi_{n})
\end{align*}
且
\begin{align*}
0\leqslant\int_{E}f-\int_{E}\varphi_{n}=\int_{E}(f-\varphi_{n})\leqslant\int_{E}(\psi_{n}-\varphi_{n}).
\end{align*}
因此
\begin{align*}
\lim_{n\to\infty}\int_{E}\varphi_{n}=\int_{E}f=\lim_{n\to\infty}\int_{E}\psi_{n}.
\end{align*}
\end{proof}

\begin{theorem}\label{theorem:ch05-有界可积可测}
令 $f$ 为有限测度集 $E$ 上的有界函数. 则 $f$ 在 $E$ 上是 Lebesgue 可积的当且仅当它是可测的.
\end{theorem}

\begin{proof}
我们已在 4.2 节证明了有限测度集上的有界可测函数是 Lebesgue 可积的. 剩下来要证明逆命题. 假定 $f$ 是可积的. 从上 Lebesgue 积分与下 Lebesgue 积分的相等性, 我们得出存在简单函数序列 $\{\varphi_{n}\}$ 和 $\{\psi_{n}\}$ 使得
\begin{align*}
\text{在 }E\text{ 上对所有 }n,\quad\varphi_{n}\leqslant f\leqslant\psi_{n}
\end{align*}
且
\begin{align*}
\lim_{n\to\infty}\int_{E}[\psi_{n}-\varphi_{n}]=0.
\end{align*}
由于一对简单函数的最大值与最小值仍然是简单的, 用积分的单调性且通过用 $\max_{1\leqslant i\leqslant n}\varphi_{i}$ 代替 $\varphi_{n}$ 和用 $\min_{1\leqslant i\leqslant n}\psi_{i}$ 代替 $\psi_{n}$, 我们可以假定 $\{\varphi_{n}\}$ 递增而 $\{\psi_{n}\}$ 递减. 根据前一个引理, 在 $E$ 上几乎处处逐点地 $\{\varphi_{n}\}\to f$. 因此 $f$ 是可测的, 这是由于它几乎处处是可测函数序列的逐点极限.
\end{proof}

在开始考虑积分的时候, 我们证明了若闭有界区间 $[a,b]$ 上的有界函数在 $[a,b]$ 上是 Riemann 可积的, 则它在 $[a,b]$ 上是 Lebesgue 可积的且两个积分相等. 因此从前一个定理我们可以推出, 若 $[a,b]$ 上的有界函数是 Riemann 可积的, 则它是可测的. 以下定理更为精确.

\begin{theorem}[Lebesgue]\label{theorem:ch05-Riemann刻画}
令 $f$ 为闭有界区间 $[a,b]$ 上的有界函数. 则 $f$ 在 $[a,b]$ 上是 Riemann 可积的当且仅当 $[a,b]$ 中使得 $f$ 不连续的点的集合的测度为零.
\end{theorem}

\begin{proof}
我们首先假定 $f$ 是 Riemann 可积的. 从 $[a,b]$ 上的上 Riemann 积分与下 Riemann 积分的相等性推出, 存在 $[a,b]$ 的分划序列 $\{P_{n}\}$ 和 $\{P_{n}'\}$ 使得
\begin{align*}
\lim_{n\to\infty}[U(f,P_{n})-L(f,P_{n}')]=0,
\end{align*}
其中 $U(f,P_{n})$ 和 $L(f,P_{n}')$ 是上与下 Darboux 和. 在加细下, 由于下 Darboux 和增加而上 Darboux 和减少, 通过用公共的加细 $P_{1},\cdots,P_{n},P_{1}',\cdots,P_{n}'$ 代替每个 $P_{n}$, 我们可以假设每个 $P_{n+1}$ 是 $P_{n}$ 的加细而 $P_{n}=P_{n}'$. 对每个指标 $n$, 定义 $\varphi_{n}$ 为与 $f$ 的分划 $P_{n}$ 相应的下阶梯函数, 即在 $P_{n}$ 的分点与 $f$ 一致, 且在由 $P_{n}$ 确定的开区间上取常数值, 该常数等于 $f$ 在该区间的下确界. 用类似的方式定义上阶梯函数 $\psi_{n}$. 根据 Darboux 和的定义, 对所有 $n$,
\begin{align*}
L(f,P_{n})=\int_{a}^{b}\varphi_{n}\quad\text{且}\quad U(f,P_{n})=\int_{a}^{b}\psi_{n}.
\end{align*}
则 $\{\varphi_{n}\}$ 和 $\{\psi_{n}\}$ 是可积函数序列, 使得对每个指标 $n$, 在 $E$ 上, $\varphi_{n}\leqslant f\leqslant\psi_{n}$. 此外, 序列 $\{\varphi_{n}\}$ 是递增的而 $\{\psi_{n}\}$ 是递减的, 这是因为每个 $P_{n+1}$ 是 $P_{n}$ 的一个加细. 最后,
\begin{align*}
\lim_{n\to\infty}\int_{a}^{b}[\psi_{n}-\varphi_{n}]=\lim_{n\to\infty}[U(f,P_{n})-L(f,P_{n})]=0.
\end{align*}
我们从前一个引理推出
\begin{align*}
\text{在 }[a,b]\text{ 上 a.e. 逐点 }\{\varphi_{n}\}\to f\text{ 且 }\{\psi_{n}\}\to f.
\end{align*}
那些使得 $\{\psi_{n}(x)\}$ 或 $\{\varphi_{n}(x)\}$ 不收敛于 $f(x)$ 的点 $x$ 的集合 $E$ 的测度为零. 令 $E_{0}$ 为 $E$ 与这些 $P_{n}$ 的所有分点的集合的并. 作为零测度集与可数集的并集, $m(E_{0})=0$. 我们宣称 $f$ 在 $E\sim E_{0}$ 的每个点连续. 事实上, 令 $x_{0}$ 属于 $E\sim E_{0}$. 为证明 $f$ 在 $x_{0}$ 连续, 令 $\varepsilon>0$. 由于 $\{\psi_{n}(x_{0})\}$ 和 $\{\varphi_{n}(x_{0})\}$ 收敛于 $f(x_{0})$, 我们可以选取自然数 $n_{0}$ 使得
\begin{align}
f(x_{0})-\varepsilon<\varphi_{n_{0}}(x_{0})\leqslant f(x_{0})\leqslant\psi_{n_{0}}(x_{0})<f(x_{0})+\varepsilon.\label{align:ch05-连续性估计}
\end{align}
由于 $x_{0}$ 不是 $P_{n_{0}}$ 的分点, 可以选取 $\delta>0$ 使得开区间 $(x_{0}-\delta,x_{0}+\delta)$ 包含在由 $P_{n_{0}}$ 确定的包含 $x_{0}$ 的开区间 $I_{n_{0}}$ 中. 这个包含蕴涵
\begin{align*}
\text{若 }|x-x_{0}|<\delta,\quad\text{则 }\varphi_{n_{0}}(x_{0})\leqslant\varphi_{n_{0}}(x)\leqslant f(x)\leqslant\psi_{n_{0}}(x).
\end{align*}
从这个不等式和不等式\eqref{align:ch05-连续性估计}我们推出
\begin{align*}
\text{若 }|x-x_{0}|<\delta,\quad\text{则 }|f(x)-f(x_{0})|<\varepsilon.
\end{align*}
因此 $f$ 在 $x_{0}$ 连续.

剩下来要证明逆命题. 假设 $f$ 在 $[a,b]$ 中的几乎所有点连续. 令 $\{P_{n}\}$ 为 $[a,b]$ 的任意分划序列使得
\begin{align*}
\lim_{n\to\infty}\operatorname{gap}P_{n}=0.
\end{align*}
我们宣称
\begin{align}
\lim_{n\to\infty}[U(f,P_{n})-L(f,P_{n})]=0.\label{align:ch05-Darboux差}
\end{align}
若这得到验证, 则从关于下 Riemann 积分和上 Riemann 积分的如下估计:
\begin{align*}
\text{对所有 }n,\quad 0\leqslant\overline{\int_{a}^{b}}f-\underline{\int_{a}^{b}}f\leqslant[U(f,P_{n})-L(f,P_{n})]
\end{align*}
我们得出 $f$ 在 $[a,b]$ 上是可积的. 对每个 $n$, 令 $\varphi_{n}$ 和 $\psi_{n}$ 为与 $f$ 在分划 $P_{n}$ 下对应的下和上阶梯函数. 要证明\eqref{align:ch05-Darboux差}即要证明
\begin{align}
\lim_{n\to\infty}\int_{a}^{b}[\psi_{n}-\varphi_{n}]=0.\label{align:ch05-阶梯积分差}
\end{align}
阶梯函数的 Riemann 积分等于它的 Lebesgue 积分. 此外, 由于函数 $f$ 在有界集 $[a,b]$ 上有界, 序列 $\{\varphi_{n}\}$ 和 $\{\psi_{n}\}$ 在 $[a,b]$ 上一致有界. 因此, 根据有界收敛定理, 为证明\eqref{align:ch05-阶梯积分差}仅须证明在 $(a,b)$ 中那些使得 $f$ 连续的点与那些不是任何分划 $P_{n}$ 的分点构成的集合上逐点地 $\{\varphi_{n}\}\to f$ 且 $\{\psi_{n}\}\to f$. 令 $x_{0}$ 为这样的点. 我们证明
\begin{align}
\lim_{n\to\infty}\varphi_{n}(x_{0})=f(x_{0})\quad\text{且}\quad\lim_{n\to\infty}\psi_{n}(x_{0})=f(x_{0}).\label{align:ch05-逐点阶梯}
\end{align}
令 $\varepsilon>0$. 令 $\delta>0$, 使得
\begin{align}
\text{若 }|x-x_{0}|<\delta,\quad f(x_{0})-\varepsilon/2<f(x)<f(x_{0})+\varepsilon/2.\label{align:ch05-局部振荡}
\end{align}
选取指标 $N$ 使得若 $n\geqslant N$, $\operatorname{gap}P_{n}<\delta$. 若 $n\geqslant N$ 而 $I_{n}$ 是由 $P_{n}$ 确定的开分划区间, 它包含 $x_{0}$, 则 $I_{n}\subseteq(x_{0}-\delta,x_{0}+\delta)$. 我们从\eqref{align:ch05-局部振荡}推出
\begin{align*}
f(x_{0})-\varepsilon/2\leqslant\varphi_{n}(x_{0})<f(x_{0})<\psi_{n}(x_{0})\leqslant f(x_{0})+\varepsilon/2.
\end{align*}
因此对所有 $n\geqslant N$,
\begin{align*}
0\leqslant\psi_{n}(x_{0})-f(x_{0})<\varepsilon\quad\text{且}\quad 0\leqslant f(x_{0})-\varphi_{n}(x_{0})<\varepsilon.
\end{align*}
因此\eqref{align:ch05-逐点阶梯}成立, 从而定理证明完毕.
\end{proof}

\begin{exercise}\label{exercise:ch05-15}
令 $f$ 和 $g$ 为在 $[a,b]$ 上 Riemann 可积的有界函数. 证明乘积 $fg$ 在 $[a,b]$ 上也是 Riemann 可积的.
\end{exercise}

\begin{exercise}\label{exercise:ch05-16}
令 $f$ 为 $[a,b]$ 上的有界函数, 其不连续点的集合的测度为零. 证明 $f$ 是可测的. 接着证明无须有界性的假设, 结论同样成立.
\end{exercise}

\begin{exercise}\label{exercise:ch05-17}
令 $f$ 为 $[0,1]$ 上的在 $(0,1]$ 上连续的函数. 证明序列 $\left\{\int_{[1/n,1]}f\right\}$ 收敛, 然而 $f$ 在 $[0,1]$ 上有可能不是 Lebesgue 可积的. 若 $f$ 是非负的这能发生吗?
\end{exercise}

\end{document}